%! AP Statistics — Unit 7, Lesson 7.2 — Student Packet Cover Sheet
\documentclass[10pt]{article}
\usepackage{apstats-article}
\usepackage{apstats-boxes}
\usepackage{tikz}

\begin{document}

\begin{tikzpicture}[remember picture, overlay]
  \fill[navy] (current page.north west) rectangle
    ([yshift=-2.2cm]current page.north east);
  \node[anchor=west, white, font=\large\bfseries] at
    ([xshift=1cm, yshift=-0.7cm]current page.north west)
    {\CourseName: \SchoolYear};
  \node[anchor=west, white, font=\normalsize] at
    ([xshift=1cm, yshift=-1.3cm]current page.north west)
    {Unit 7: Inference for Quantitative Data: Means};
  \node[anchor=west, white, font=\normalsize\bfseries] at
    ([xshift=1cm, yshift=-1.9cm]current page.north west)
    {Lesson 7.2 \quad Constructing a Confidence Interval for a Population Mean};
\end{tikzpicture}

\vspace{2.4cm}
\namedateperiod

\begin{learningtargetbox}
\begin{itemize}
    \item I can describe the $t$-distribution and explain how it differs from the
          standard normal distribution.
    \item I can identify the one-sample $t$-interval as the appropriate procedure for
          estimating a population mean when $\sigma$ is unknown, and verify the required
          conditions (Random, 10\%, Large Sample).
    \item I can calculate a one-sample $t$-interval using $\bar x \pm t^* \cdot s/\sqrt{n}$.
    \item I can apply the one-sample $t$-interval to a matched-pairs design by analyzing
          the distribution of differences.
\end{itemize}
\end{learningtargetbox}

\begin{tocbox}
\renewcommand{\arraystretch}{1.5}
\begin{tabularx}{\linewidth}{clXc}
    \textbf{\#} & \textbf{Component} & \textbf{Description} & \textbf{Score} \\
    \hline
    1 & Warm-Up      & Spiral review: $z$-interval for proportions; intro to $t$ & \hfill/6 \\
    2 & Guided Notes & $t$-distribution, conditions, formula, matched pairs         & \hfill/12 \\
    3 & Activity     & Partner practice: full four-step $t$-interval procedure       & \hfill/15 \\
    4 & Exit Ticket  & Independent: construct and interpret a $t$-interval           & \hfill/6 \\
    5 & Homework     & Practice: $t$-intervals and matched pairs                     & \hfill/15 \\
    \hline
       & \multicolumn{2}{l}{\textbf{Total}} & \hfill/54 \\
\end{tabularx}
\end{tocbox}

\end{document}
